\section{Introduction}\label{sec:intro}

On 11 June 2026 eight authors, writing in a personal capacity, published \emph{Europe 2031}, a scenario of what getting artificial intelligence wrong would mean for Europe \citep{juijn2026europe}. They call ASML, the Dutch firm that alone builds extreme ultraviolet (EUV) lithography machines, ``the only card Europe still holds''. In the story's climax, in March 2031, the White House proposes to the leaders of the Netherlands, Germany and France that ASML be folded into a joint Dutch-American holding company in which Washington holds a controlling vote over output, customers and technology transfer. When the three leaders decline, Washington makes the offer public and threatens to drop the whole EU to Tier~3 of its rules on access to frontier AI. Three years earlier in the story, the Dutch had ended ASML's remaining sales and servicing of immersion tools in China under the threat of the US foreign direct product rule, and one character reflects that ``we didn't even manage to negotiate a single thing in return''.

The scenario is fiction (\cref{app:scenario} sets its episodes beside the model), but its premise describes the record since 2019 (\cref{sec:background}). In 2019, after US lobbying reported by Reuters, the Dutch government let ASML's licence to ship an EUV machine to China lapse without renewal \citep{nltimes2020pressured}. National licensing of advanced immersion lithography took effect on 1 September 2023 \citep{govnl2023export}. In September 2024 licensing for the older NXT:1970i and NXT:1980i moved from Washington to The Hague \citep{asml2024license}, a measure on metrology and inspection systems published in January 2025 took effect on 1 April 2025, and in November 2025 the EU put the most advanced immersion tools on the Union control list through Delegated Regulation 2025/2003 \citep{asml2026annual}. What The Hague received was a status. In December 2024 the Bureau of Industry and Security (BIS) placed the Netherlands on Supplement No.~4, whose listed countries' exporters are excluded from new foreign direct product and de minimis requirements on chipmaking equipment, because ``Governments in those destinations \ldots\ have the authority to control key SME items of concern'' \citep{bis2024fdpr}. BIS can withdraw the status by amending a rule. We code it as a threat left unexercised for as long as Europe controls the items itself, a reward for self-binding and not a transfer (\cref{sec:coding}).

Yet the holders of these chokepoints resist some demands. In April 2026 bills in both chambers of Congress, the MATCH Act, proposed controls on all immersion lithography, a policy of denial for licences to service covered Chinese fabs, and extended US controls on allies whose own controls were not certified as equivalent within 150 days \citep{hr8170, s4281}. Prime Minister Jetten raised the proposed restrictions with President Trump at a White House dinner on 13 April \citep{nltimes2026pushback}, and on 12 May the trade minister told parliament that ``the cabinet is against the extraterritorial effect of the American proposal'' \citep[our translation]{tk2026match}. As of June 2026 both the Netherlands and Japan still let their firms service equipment in Chinese fabs \citep{zakaria2026match}. Tokyo had put 23 types of chipmaking equipment under licence in 2023 while stating that it was ``not aligning itself with the United States' measures'' \citep{allen2023translation, hoganlovells2023japan}. Germany reportedly resisted US, Dutch and Japanese pressure to join controls aimed above all at Zeiss, ASML's sole supplier of optics; in one European executive's words, ``Germany does not want to join this club'' \citep{barkin2024zeiss}. South Korea, by contrast, accepted the loss of validated end-user status for the Chinese fabs of Samsung and SK hynix in 2025 \citep{bis2025veu}.

Holders that Washington needs comply with little compensation, and they comply with some demands while resisting others. Each of the accounts that should explain this pattern misses part of it. Theories of asymmetric and weaponized interdependence give power to the less dependent party and to the state with jurisdiction over a hub \citep{hirschman1945national, keohane1977power, farrell2019weaponized}. On that account Washington, which has no substitute for EUV lithography, should be paying The Hague. Accounts of alliance dependence and hierarchy treat compliance as the price of protection already received \citep{morrow1991alliances, lake2009hierarchy}, and \citet{drezner1999sanctions} finds that coercion among allies, though rare, often wins concessions. They fit the compliance but leave the resistance on servicing and extraterritorial reach unexplained. Models of competition among bidders hold that a client courted by two patrons is paid what the losing bidder would offer \citep{groseclose1996buying, bueno2016competition}, as Kyrgyzstan was when it used a Russian aid package to raise the rent Washington paid for the Manas base in 2009 \citep{cooley2012great}. Customers in China took 36 percent of ASML's sales in 2024 \citep{asml2026annual}. A bidding account predicts a price, and The Hague has not been paid one.

We argue that a chokepoint holder's leverage is bounded by its exit premium, and that a patron's demands are bounded by the holder's slack. The exit premium is the rival's credible offer, net of the costs the patron can credibly impose on exit and the value the holder places on alignment; slack is the amount by which the premium falls short of zero. The first bound is a known result: the outside-option principle \citep{shaked1984involuntary, sutton1986noncooperative, binmore1989outside} and the one-voter case of sequential vote buying \citep{groseclose1996buying, dekel2008vote}. Its substance goes back to \citet{hirschman1945national} and \citet{wagner1988economic}, and \citet{clayton2026framework} derive a two-hegemon version. The second bound organizes the paper. \citet{clayton2026putting} write that ``threats generate power to the extent that they generate slack in the participation constraint''. We take their term and ask what slack buys a patron when each of its threats must be sequentially rational, when the holder's outside option is a rival's bid, and when a threat that is carried out is spent. Our answer is a slack principle: whatever instruments back them, the conditional parts of a patron's threats draw on one budget, and every demand the patron makes (on sales, servicing, hedging or control) competes for it. A holder with a negative premium therefore gets nothing for its chokepoint unless the dispute risks rupture, however much its patron needs it. Its patron must ration demands. A demand meets resistance when its cost to the holder exceeds the bite that remains, and China can raise that cost. Its 2026 decree on supply-chain security lets Beijing restrict the trade and investment in China of foreign firms that cut off normal trade with Chinese entities \citep{statecouncil2026supplychain}, which on our reading attaches a counter-threat to any order to stop servicing Chinese fabs.

In the game of \cref{sec:model}, the United States and China bid for Europe's alignment and one commitment technology governs every threat and promise. After the bidding benchmark (\cref{lem:bidding}), \cref{prop:master} makes China's offer Europe's outside option in random-proposer bargaining with a hazard of rupture. The transfer from Washington is $P^*=\max\{Q,\psi(Q+\beta(V-Q))\}$, where $Q$ is the exit premium, $V$ Washington's dependence, $\beta$ Europe's probability of proposing and $\psi$ the rupture weight. \Cref{cor:master} states the threshold and names the five channels through which $V$ enters. \Cref{prop:slack} proves the slack principle: in every subgame perfect equilibrium of any sequence of ultimatums, Europe's total loss is bounded by the bite of the credible threats, the unconditional parts plus the smaller of the conditional parts and the slack. \Cref{cor:rationing} turns the bound into a rationing problem whose capacity shrinks as the premium rises toward zero. By \cref{prop:counter}, China executes a counter-threat only to crowd out another US demand, and when Washington has a divisible demand China cannot threaten, the counter-threat taxes the patron and leaves Europe's loss unchanged. \Cref{prop:conditionality} makes conditionality a choice with an assurance cost \citep{pauly2024damned}. A patron makes its guarantee more conditional than deterrence requires only to create slack it can spend, so conditionality is the price of concessions, and a stronger Chinese bid can be met with a more conditional guarantee instead of a payment.

\Cref{sec:control} shows that a wasting slack gives Washington a preventive motive to demand control of the chokepoint while its threats still bite. \Cref{prop:control} weighs the rents a transfer removes and Washington's intrinsic value of control against the issuing cost, Europe's sovereignty cost and China's punishment of the transfer. By \cref{cor:control}, an autonomy investment that Europe cannot commit to forgo strengthens the motive, and a credible Chinese threat against a transfer can protect European ownership without being carried out. \Cref{cor:windows,rem:bind,prop:breakdown} cover the timing of demands, self-binding and refusal. \Cref{sec:decide} drops the unitary holder (\cref{rem:gatekeeper,rem:firm,rem:council}) and gives Washington's dependence its most direct role: interrupting servicing during a dispute helps Europe only if the US buffer of tools and know-how is short and only for the servicing Washington cannot compel (\cref{prop:pain,cor:interrupt}).

\Cref{sec:evidence} is a plausibility probe. It derives predictions that separate the slack logic from the rival accounts (\cref{tab:implications}), reads the Dutch record and the Nexperia dispute of 2025 through the model (\cref{sec:record}), and compares holders in Japan, South Korea, Germany and Taiwan (\cref{sec:cases}). The record is consistent with the jurisdictional channel: Washington has coerced mainly by granting and withdrawing statuses it can change by rule. It leaves the security channel largely untested. We found no public record that links European chip controls to US force posture, so the evidence cannot show how conditional the guarantee is on export controls, and \cref{sec:tests} describes the designs that could.

We contribute to four literatures.
\begin{enumerate}[label=\arabic*.]
\item \emph{Geoeconomics and coercion.} In \citet{clayton2026framework}, pooled threats are enforcement threats that raise the target's inside option and are never spent, and threats to cut a target off appear only in their Supplement B.3.1. The bidding models cited above have carrots only, and \citet{lacy2004theory} model linkage through preferences and information. We put a rival's bid inside the participation constraint, split each instrument into conditional and unconditional parts, and let executed threats be spent. The result is a theory of how a patron rations coercion across demands as an ally's outside option improves (\cref{prop:slack,cor:rationing,prop:counter}).
\item \emph{Alliance politics.} Alliance theory already separates an ally's dependence from the patron's commitment \citep{snyder1984security, snyder1997alliance}, and abandonment threats are known to work when they are conditional and paired with assurance \citep{gerzhoy2015alliance, blankenship2023burden}. We make conditionality a costly choice and show that its value to the patron is the slack it creates against the rival's credible offer (\cref{prop:conditionality}). Base politics has the same structure with territory as the asset \citep{cooley2008base}; the ASML case adds nested jurisdiction, hold-up through servicing and preventive control.
\item \emph{Shifting power and control rights.} Preventive war answers a rising state's inability to commit \citep{fearon1995rationalist, powell2006war, krainin2017preventive}, states try to stop investments in future bargaining power \citep{debs2014known, bas2016dynamic, benson2023commitment}, and control rights are the remedy for hold-up \citep{grossman1986costs, lake1996anarchy, cooley2009contracting}. We show that a patron's power to demand control is capped by the holder's current slack, so the demand comes in windows while that slack lasts (\cref{prop:control,cor:windows}). When Washington places little intrinsic value on control, an accepted transfer is inefficient in the way preventive war is: Europe bears its sovereignty cost so that Washington avoids rents that are only transfers.
\item \emph{Weaponized interdependence.} \citet{farrell2019weaponized} make jurisdiction over a hub and the domestic institutions to use it the conditions for weaponizing it. When the holder of a hub bargains with a state that depends on it, we add a third: the holder's exit premium must be positive.
\end{enumerate}

\Cref{sec:background} describes ASML's position and the control sequence, and \cref{sec:literature} the work we build on. \Cref{sec:model,sec:slack,sec:control,sec:decide} present the model and results. \Cref{sec:evidence} reads the record and the comparison cases through the model, and \cref{sec:implications} draws implications and states the scope conditions (\cref{sec:scope}). Proofs are in \cref{app:proofs} and numerical checks in \cref{app:numerics}.
